% §17.5 Boolean Algebra Over Trapdoor Sets
%
% Cross-reference resolution status:
%
% DEFINED here:
%   \label{sec:boolean-trapdoor}              -- this section
%   \label{def:trapdoor-set}                  -- trapdoor set definition
%
% ALREADY DEFINED (chapter 3):
%   \Cref{thm:composition}                    -- chapter 3 §3.3 composition theorem
%
% ALREADY DEFINED (chapter 5):
%   \Cref{sec:bloom-from-scratch}             -- chapter 5 §5.2 Bloom filter
%
% ALREADY DEFINED (chapter 6):
%   \Cref{sec:cancellation}                   -- chapter 6 §6.3 cancellation analysis
%   \Cref{def:perfect-filter-line}            -- chapter 6 §6.4 definition
%   \Cref{prop:complement-axis-swap}          -- chapter 6 §6.3 complement axis swap
%
% INTERNAL refs (chapter 17):
%   \Cref{def:cipher-map}                     -- §17.1
%   \Cref{sec:trapdoor-paradigm}              -- §17.3
%   \Cref{thm:cipher-composition}             -- §17.2
%
% §17.5 introduces Boolean algebra over trapdoor sets. Five items:
% (1) trapdoor sets: encrypted set representations needing key for
% membership; Bloom is example; (2) Boolean ops (union, intersection,
% complement) on trapdoor representations; (3) composition theorem
% (chapter 3) applies to trapdoor-set composition; (4) cancellation
% (chapter 6 §6.3) lets trapdoor composition escape perfect-filter
% line; (5) worked example: query x in A AND x in B on two trapdoor
% sets, derive overall fprate.

\section{Boolean Algebra Over Trapdoor Sets}
\label{sec:boolean-trapdoor}

The trapdoor cipher maps of \Cref{sec:trapdoor-paradigm} are
useful in isolation: a single encrypted Bloom filter answers
membership queries with controlled error. They become more useful
when several are combined. A common request in encrypted-search
applications is the conjunctive query: \emph{is $x$ in $A$ AND in
$B$?}, asked of two trapdoor sets held on the same server. The
answer requires Boolean reasoning over encrypted representations,
not over plaintext sets. This section asks how Boolean algebra
behaves when its operands are trapdoor sets, and finds that the
framework's composition rules continue to apply, with the same
geometric subtleties that chapter 6 raised for plaintext approximate
sets.

\subsection*{Trapdoor sets}

A trapdoor set is an encrypted representation of a plaintext set
whose membership query requires a secret key.

\begin{definition}[Trapdoor set]
\label{def:trapdoor-set}
Let $A \subseteq X$ be a plaintext set and let $k$ be a secret key.
A \emph{trapdoor set} for $A$ under $k$, denoted $\tilde A$, is an
encrypted data structure together with an associated membership
cipher map
\[
  \tilde m_A : \mathsf{Enc}_k(X) \to \mathsf{Enc}_k(\{0, 1\})
\]
such that, for $x \in X$, $\mathsf{Dec}_k(\tilde m_A(\mathsf{Enc}_k(x))) =
\mathbf{1}[x \in A]$ with bounded error rates $\fprate_A, \fnrate_A$ in
the sense of \Cref{def:cipher-map}. Without the key $k$, the
membership query is computationally infeasible; with the key, it is
as cheap as a plaintext set lookup.
\end{definition}

The Bloom filter of chapter 5 \Cref{sec:bloom-from-scratch} is the
canonical example. A server holds an encrypted bit array
representing the Bloom filter's state, plus encrypted hash family
parameters that depend on the key. The client encrypts a query
element under the key, sends the ciphertext to the server, and the
server returns the encrypted membership bit. The cipher map has
$\fprate_A = (1 - e^{-kn/m})^k$ (the standard Bloom-filter formula
with $k$ hash functions, $n$ inserted elements, and $m$ bits), and
$\fnrate_A = 0$ from the one-sided error of the construction. The
trapdoor set sits at the point $(\fprate_A, 0)$ on the perfect-filter
line of \Cref{def:perfect-filter-line}.

Other trapdoor-set constructions sit at different points on the
rate plane. A perfect-hash-filter trapdoor set (chapter 15) sits
at $(2^{-r}, 0)$ for an $r$-bit fingerprint. A counting trapdoor
set built on the same hash family can support deletions at the
cost of a small $\fnrate_A$. The space of constructions is large;
what they share is the structural property of \Cref{def:trapdoor-set}:
encrypted representation, key-dependent membership query, bounded
Bernoulli error.

\subsection*{Boolean operations on trapdoor sets}

The three primitive Boolean operations on sets are union $A \cup B$,
intersection $A \cap B$, and complement $\complement A$.
Each lifts to the trapdoor setting, but the lift requires the server
to compute on encrypted representations rather than plaintext.

\textbf{Union.} A trapdoor set $\tilde{A \cup B}$ realizes the
plaintext union $A \cup B$ in the sense that
\[
  \mathsf{Dec}_k(\tilde m_{A \cup B}(\mathsf{Enc}_k(x))) =
  \mathbf{1}[x \in A \cup B] \;=\; \mathbf{1}[x \in A] \lor \mathbf{1}[x \in B],
\]
up to Bernoulli error. For Bloom-filter trapdoor sets the union can
be computed by a bitwise OR of the encrypted bit arrays, provided
the encryption scheme is additively homomorphic over the bits. The
server returns the encrypted bit obtained by querying both arrays
and OR-ing the results, and the client decrypts to read the union
membership.

\textbf{Intersection.} The intersection $\tilde{A \cap B}$ is the
analogous lift of $A \cap B$. The server queries both trapdoor sets
and AND-s the encrypted membership bits, returning an encrypted
result that decrypts to $\mathbf{1}[x \in A] \land \mathbf{1}[x \in B]$.
Unlike union, intersection on Bloom-filter trapdoor sets cannot
generally be computed by a bitwise operation on the underlying
arrays (the bitwise AND of two Bloom-filter arrays is not the Bloom
filter of the intersection, a well-known pitfall in the cleartext
case). Instead the server runs both membership queries separately
and combines the encrypted output bits with a homomorphic AND.

\textbf{Complement.} The complement $\tilde{\complement A}$ is
computed by complementing the encrypted membership bit:
$\mathsf{Dec}_k(\tilde m_{\complement A}(\mathsf{Enc}_k(x))) =
1 - \mathbf{1}[x \in A]$ up to Bernoulli error. The server flips
the returned bit (or runs the cipher map in complement mode). The
operation is cheap; what changes, as \Cref{prop:complement-axis-swap}
warned, is the position of the result on the perfect-filter line.

These three primitives suffice to build any finite Boolean
combination by composition. The interesting questions concern how
the operations' error rates compose, and how confidentiality
propagates when the server has access to the operations' encrypted
outputs.

\subsection*{The composition theorem applies}

The error rates of a Boolean combination of trapdoor sets are the
error rates of the corresponding plaintext Boolean operation,
computed by the framework's composition theorem.
\Cref{thm:composition} stated the rule for plaintext approximate
sets: when an operation combines $k$ approximate operands whose
error events are mutually independent under Axiom 1, the resulting
overall error rate satisfies
\[
  \eta_{\text{total}} \;=\; 1 - \prod_{i=1}^{k} (1 - \eta_i),
\]
applied separately to the FPR and FNR channels.

The same formula applies to trapdoor-set composition without
modification. Two trapdoor sets $\tilde A, \tilde B$ with per-set
rates $(\fprate_A, 0)$ and $(\fprate_B, 0)$, combined under union,
yield a trapdoor set $\tilde{A \cup B}$ with rates
\[
  \fprate_{A \cup B} \;=\; 1 - (1 - \fprate_A)(1 - \fprate_B),
  \qquad
  \fnrate_{A \cup B} \;=\; 0.
\]
Combined under intersection the rates are
\[
  \fprate_{A \cap B} \;=\; \fprate_A \cdot \fprate_B,
  \qquad
  \fnrate_{A \cap B} \;=\; 0,
\]
because a query $x \notin A \cap B$ is a false positive only if it
is a false positive under both component membership tests, which
happens with probability $\fprate_A \cdot \fprate_B$ under
independence.

The conditional-independence hypothesis is the chapter 4 Axiom 1
applied to the trapdoor setting: the two trapdoor sets' cipher
maps use independent randomness (independent hash family seeds,
independent encryption noise) and so produce independent false-positive
events for any given query. This is automatic if the two trapdoor
keys are generated independently, and is the standard assumption in
encrypted-search constructions.

The same algebra propagates through deeper compositions. A
three-way conjunctive query $x \in A \cap B \cap C$ on trapdoor
sets at rates $\fprate_A, \fprate_B, \fprate_C$ has overall
$\fprate = \fprate_A \fprate_B \fprate_C$. A $k$-way conjunction
has $\fprate = \prod_i \fprate_i$, which decays exponentially in
$k$: deep conjunctive queries are highly precise even when the
individual trapdoor sets are loose. The framework's composition
law captures this strengthening exactly, and the trapdoor-set
geometry inherits it.

\subsection*{Cancellation escapes the perfect-filter line}

The composition formula above kept the operands on the perfect-filter
line: both unions and intersections of FNR-zero trapdoor sets stayed
FNR-zero. What chapter 6 \Cref{sec:cancellation} flagged is that
this clean behavior is not a general property of Boolean composition;
it depends on the operands sitting on the same axis. As soon as the
complement enters the picture (or as soon as one operand sits on the
opposite axis from another), the result can leave the perfect-filter
line entirely.

The cancellation phenomenon transfers directly to trapdoor sets.
Consider a trapdoor union $\tilde A \cup \tilde B$ at rates
$(\fprate_A, 0)$ and $(\fprate_B, 0)$, both on the FNR-zero axis.
Now take the complement of this union, $\tilde{\complement(A \cup B)}$.
By \Cref{prop:complement-axis-swap} the complement of an FNR-zero
trapdoor set is FPR-zero, so $\tilde{\complement(A \cup B)}$ sits
at the point $(0, \fprate_{A \cup B})$ on the FPR-zero axis. The
trapdoor system has moved off the original axis.

Compose further. Suppose the server is asked $\tilde{\complement(A
\cup B)} \cap \tilde C$, where $\tilde C$ is a third trapdoor set
back on the FNR-zero axis at rates $(\fprate_C, 0)$. The
intersection now mixes operands from different axes. By the rate
algebra,
\[
  \fprate_{\complement(A \cup B) \cap C} \;=\;
  \fprate_{\complement(A \cup B)} \cdot \fprate_C \;=\; 0,
\]
because one operand is FPR-zero. But the FNR is no longer zero,
because the FPR-zero operand carries a positive $\fnrate$:
\[
  \fnrate_{\complement(A \cup B) \cap C} \;=\;
  1 - (1 - \fprate_{A \cup B})(1 - \fnrate_C)
  \;=\; \fprate_{A \cup B},
\]
since $\fnrate_C = 0$. The composed trapdoor set has rates
$(0, \fprate_{A \cup B})$ and now sits on the FPR-zero axis: the
two complements and the cross-axis intersection have moved the
result through the rate plane in a manner the perfect-filter-line
classification cannot describe by a single one-sidedness flag.

More elaborate Boolean combinations can produce trapdoor sets with
both rates positive, fully off the perfect-filter line. This is
the same cancellation phenomenon chapter 6 documented for
plaintext approximate sets, lifted intact to the trapdoor setting.
The practical implication is that designing trapdoor-set Boolean
combinations requires the same care as designing plaintext Boolean
combinations: track the operand axes through the composition, and
expect the perfect-filter-line classification to dissolve as soon
as the composition mixes axes.

\subsection*{Worked example: conjunctive query}

A concrete example fixes the rate algebra. Take two trapdoor Bloom
filters $\tilde A$ and $\tilde B$ on plaintext sets $A, B \subseteq
X$, each configured at $\fprate = 0.01$ with $\fnrate = 0$. The
client wants to answer the query
\[
  q(x) \;=\; \mathbf{1}[x \in A] \;\land\; \mathbf{1}[x \in B] \;=\;
  \mathbf{1}[x \in A \cap B].
\]
The client encrypts $x$ under the shared key, sends $\mathsf{Enc}_k(x)$
to the server, and asks the server to compute the conjunctive
membership cipher map.

\textbf{Server computation.} The server queries both trapdoor sets,
producing encrypted bits $\tilde b_A = \tilde m_A(\mathsf{Enc}_k(x))$
and $\tilde b_B = \tilde m_B(\mathsf{Enc}_k(x))$. It combines them
under a homomorphic AND to produce $\tilde b = \tilde b_A \land
\tilde b_B$, and returns $\tilde b$ to the client. The client
decrypts: $\mathsf{Dec}_k(\tilde b)$ is a bit in $\{0, 1\}$ which
approximates $\mathbf{1}[x \in A \cap B]$ up to the composed
Bernoulli error.

\textbf{Rate analysis.} Apply the composition formula above. The
intersection of two FNR-zero trapdoor sets has $\fnrate = 0$ (no
real positive is lost, because both component queries return $1$
on true positives). The composed FPR is the product:
\[
  \fprate_{A \cap B} \;=\; \fprate_A \cdot \fprate_B \;=\; 0.01 \times 0.01
  \;=\; 10^{-4}.
\]
The conjunctive query has $\fprate = 10^{-4}$, two orders of
magnitude tighter than either component trapdoor set. The client
receives an answer that is correct on all true positives (because
$\fnrate = 0$) and rejects all but one in ten thousand false
positives.

\textbf{Interpretation.} The intersection sharpens precision
multiplicatively. Each additional trapdoor set in a $k$-way
conjunction divides the false-positive rate by another factor on
the order of $\fprate$, and a $k$-way conjunction has $\fprate
= \fprate^k$ (assuming equal-rate operands and independent
trapdoors). The framework's composition theorem captures this
sharpening exactly; the trapdoor setting inherits it without
qualification.

This is the most useful pattern in encrypted search built from
trapdoor sets. A query that asks for documents matching all of
several encrypted keywords becomes a conjunction of trapdoor-set
memberships, and the conjunction's precision sharpens exponentially
in the number of keywords. The same composition law that governs
plaintext Bloom-filter conjunctions governs the trapdoor lift:
matrix multiplication in the channel-matrix formulation of
\Cref{thm:cipher-composition}, or equivalently the product
formula above when the operands are FNR-zero. The next section turns
to how this rate-algebra story is also a confidentiality-algebra
story, and asks how much information about the plaintext sets the
server can recover from the conjunctive query stream.
